% GENERATED FILE. Edit canonical lessons, then run npm run generate:books.
% canonical-source-hash: fnv1a64:d0e0e6b6b536adeb
% canonical-lessons: BN-C536-suparish, BN-C536-sommoti, BN-C536-rokom, BN-C536-dhoron, BN-C536-khap-khaoya

\chapter{Kinds and Choices}
\label{ch:bn-kinds-and-choices}
\hlchaptermodality{536}

\begin{chapteropening}
\textbf{By the end of this chapter:} \emph{I can talk about kinds, styles, recommendations and what fits.}

Complete the last lesson of chapter 536: I can talk about kinds, styles, recommendations and what fits.
\end{chapteropening}

\section[\texorpdfstring{\bn{সুপারিশ} (supārish)}{সুপারিশ (supārish)}]{\bn{সুপারিশ} (supārish) — a recommendation}

\label{lesson:BN-C536-suparish}



\begin{quote}
Before the new one: say the Bengali for well protected, then the Bengali for a fire extinguisher.
\end{quote}

\subsection*{You'll want to know: \bn{সুপারিশ}}
\textbf{\bn{সুপারিশ}} — \emph{supārish} — \textquotedblleft{}a recommendation\textquotedblright{}.

\begin{quote}
\bn{ডাক্তার} \bn{সুপারিশ} \bn{করলে} \bn{আমি} \bn{এই} \bn{ওষুধটা} \bn{কিনব।}
\end{quote}

\begin{quote}
\emph{Ḍāktār supārish kôrle āmi ei oshudhṭā kinbô.}
\end{quote}

\begin{quote}
If the doctor recommends it, I'll buy this medicine.
\end{quote}

\begin{cousinweb}
\textbf{\bn{সুপারিশ}} (\emph{supārish}) is usually traced to Persian \emph{sifārish}, \textquotedblleft{}a commission, a recommendation\textquotedblright{}.
\end{cousinweb}

\begin{grammarlens}[title={What you've built}]
One more word for saying what happens if.
\end{grammarlens}

\subsection*{Your turn}
\begin{itemize}
\raggedright
  \item \emph{Say it:} \emph{supārish}
  \item \emph{Say it:} \emph{supārish}, once more
  \item \emph{Say it:} \emph{Ḍāktār supārish kôrle āmi ei oshudhṭā kinbô.} — If the doctor recommends it, I'll buy this medicine.
  \item \emph{Recall:} say the Bengali for well protected, then the Bengali for a fire extinguisher, then say \emph{supārish} again
\end{itemize}

\begin{tcolorbox}[breakable,colback=teal!4,colframe=teal!35!black,title={Before you move on}]
What does \bn{সুপারিশ} mean? (\textquotedblleft{}A recommendation\textquotedblright{}.) Say it once more.
\end{tcolorbox}
\section[\texorpdfstring{\bn{সম্মতি} (sômmôti)}{সম্মতি (sômmôti)}]{\bn{সম্মতি} (sômmôti) — consent}

\label{lesson:BN-C536-sommoti}



\begin{quote}
Before the new one: say the Bengali for a fire extinguisher, then the Bengali for a recommendation.
\end{quote}

\subsection*{You'll want to know: \bn{সম্মতি}}
\textbf{\bn{সম্মতি}} — \emph{sômmôti} — \textquotedblleft{}consent\textquotedblright{}.

\begin{quote}
\bn{বাবা} \bn{সম্মতি} \bn{দিলে} \bn{আমরা} \bn{বিয়ে} \bn{করব।}
\end{quote}

\begin{quote}
\emph{Bābā sômmôti dile āmrā biye kôrbô.}
\end{quote}

\begin{quote}
If my father gives his consent, we'll get married.
\end{quote}

\begin{cousinweb}
\textbf{\bn{সম্মতি}} (\emph{sômmôti}) is a tatsama, Sanskrit \emph{sammati}, \textquotedblleft{}sameness of opinion, approval\textquotedblright{} (Monier-Williams).
\end{cousinweb}

\begin{grammarlens}[title={What you've built}]
One more word for saying what happens if.
\end{grammarlens}

\subsection*{Your turn}
\begin{itemize}
\raggedright
  \item \emph{Say it:} \emph{sômmôti}
  \item \emph{Say it:} \emph{sômmôti}, once more
  \item \emph{Say it:} \emph{Bābā sômmôti dile āmrā biye kôrbô.} — If my father gives his consent, we'll get married.
  \item \emph{Recall:} say the Bengali for a fire extinguisher, then the Bengali for a recommendation, then say \emph{sômmôti} again
\end{itemize}

\begin{tcolorbox}[breakable,colback=teal!4,colframe=teal!35!black,title={Before you move on}]
What does \bn{সম্মতি} mean? (\textquotedblleft{}Consent\textquotedblright{}.) Say it once more.
\end{tcolorbox}
\section[\texorpdfstring{\bn{রকম} (rôkôm)}{রকম (rôkôm)}]{\bn{রকম} (rôkôm) — a kind; a sort}

\label{lesson:BN-C536-rokom}



\begin{quote}
Before the new one: say the Bengali for a recommendation, then the Bengali for consent.
\end{quote}

\subsection*{You'll want to know: \bn{রকম}}
\textbf{\bn{রকম}} — \emph{rôkôm} — \textquotedblleft{}a kind; a sort\textquotedblright{}.

\begin{quote}
\bn{দোকানে} \bn{তিন} \bn{রকম} \bn{চা} \bn{আছে।}
\end{quote}

\begin{quote}
\emph{Dokāne tin rôkôm chā āchhe.}
\end{quote}

\begin{quote}
The shop has three kinds of tea.
\end{quote}

\begin{cousinweb}
\textbf{\bn{রকম}} (\emph{rôkôm}) is usually traced to Arabic \emph{raqm}, \textquotedblleft{}a mark, a figure\textquotedblright{}, through Persian.
\end{cousinweb}

\begin{grammarlens}[title={What you've built}]
One more word for saying what happens if.
\end{grammarlens}

\subsection*{Your turn}
\begin{itemize}
\raggedright
  \item \emph{Say it:} \emph{rôkôm}
  \item \emph{Say it:} \emph{rôkôm}, once more
  \item \emph{Say it:} \emph{Dokāne tin rôkôm chā āchhe.} — The shop has three kinds of tea.
  \item \emph{Recall:} say the Bengali for a recommendation, then the Bengali for consent, then say \emph{rôkôm} again
\end{itemize}

\begin{tcolorbox}[breakable,colback=teal!4,colframe=teal!35!black,title={Before you move on}]
What does \bn{রকম} mean? (\textquotedblleft{}A kind; a sort\textquotedblright{}.) Say it once more.
\end{tcolorbox}
\section[\texorpdfstring{\bn{ধরন} (dhôrôn)}{ধরন (dhôrôn)}]{\bn{ধরন} (dhôrôn) — a manner; a style; a type}

\label{lesson:BN-C536-dhoron}



\begin{quote}
Before the new one: say the Bengali for consent, then the Bengali for a kind of thing.
\end{quote}

\subsection*{You'll want to know: \bn{ধরন}}
\textbf{\bn{ধরন}} — \emph{dhôrôn} — \textquotedblleft{}a manner; a style; a type\textquotedblright{}.

\begin{quote}
\bn{এই} \bn{বাড়ির} \bn{ধরন} \bn{পুরনো।}
\end{quote}

\begin{quote}
\emph{Ei bāṛir dhôrôn purôno.}
\end{quote}

\begin{quote}
This house is old in style.
\end{quote}

\begin{cousinweb}
\textbf{\bn{ধরন}} (\emph{dhôrôn}) is probably built on \bn{ধরা}, \textquotedblleft{}to hold\textquotedblright{} (CDIAL 6747): the way a thing is held or shaped.
\end{cousinweb}

\begin{grammarlens}[title={What you've built}]
One more word for saying what happens if.
\end{grammarlens}

\subsection*{Your turn}
\begin{itemize}
\raggedright
  \item \emph{Say it:} \emph{dhôrôn}
  \item \emph{Say it:} \emph{dhôrôn}, once more
  \item \emph{Say it:} \emph{Ei bāṛir dhôrôn purôno.} — This house is old in style.
  \item \emph{Recall:} say the Bengali for consent, then the Bengali for a kind of thing, then say \emph{dhôrôn} again
\end{itemize}

\begin{tcolorbox}[breakable,colback=teal!4,colframe=teal!35!black,title={Before you move on}]
What does \bn{ধরন} mean? (\textquotedblleft{}A manner; a style; a type\textquotedblright{}.) Say it once more.
\end{tcolorbox}
\section[\texorpdfstring{\bn{খাপ} \bn{খাওয়া} (khāp khāoyā)}{খাপ খাওয়া (khāp khāoyā)}]{\bn{খাপ} \bn{খাওয়া} (khāp khāoyā) — to fit; to suit}

\label{lesson:BN-C536-khap-khaoya}



\begin{quote}
Before the new one: say the Bengali for a kind of thing, then the Bengali for a style.
\end{quote}

\subsection*{You'll want to know: \bn{খাপ} \bn{খাওয়া}}
\textbf{\bn{খাপ} \bn{খাওয়া}} — \emph{khāp khāoyā} — \textquotedblleft{}to fit; to suit\textquotedblright{}.

\begin{quote}
\bn{এই} \bn{চাবিটা} \bn{তালায়} \bn{খাপ} \bn{খায়} \bn{না।}
\end{quote}

\begin{quote}
\emph{Ei chābiṭā tālāy khāp khāy nā.}
\end{quote}

\begin{quote}
This key doesn't fit the lock.
\end{quote}

Past (he/she …): \bn{সেই} \bn{ঢাকনা} \bn{হাঁড়িতে} \bn{খাপ} \bn{খেল}, \emph{sei ḍhāknā hā̃ṛite khāp khelô}.

\begin{cousinweb}
\textbf{\bn{খাপ} \bn{খাওয়া}} (\emph{khāp khāoyā}) joins \bn{খাপ}, \textquotedblleft{}a sheath, a fit\textquotedblright{}, to \bn{খাওয়া}, \textquotedblleft{}to eat\textquotedblright{}; CDIAL lists Bengali \emph{khāpā}, \textquotedblleft{}to become adjusted\textquotedblright{}, under a reconstructed \emph{khappa} (CDIAL 3816).
\end{cousinweb}

\begin{grammarlens}[title={What you've built}]
One more word for saying what happens if.
\end{grammarlens}

\subsection*{Your turn}
\begin{itemize}
\raggedright
  \item \emph{Say it:} \emph{khāp khāoyā}
  \item \emph{Say it:} \emph{khāp khāoyā}, once more
  \item \emph{Say it:} \emph{Ei chābiṭā tālāy khāp khāy nā.} — This key doesn't fit the lock.
  \item \emph{Recall:} say the Bengali for a kind of thing, then the Bengali for a style, then say \emph{khāp khāoyā} again
\end{itemize}

\begin{tcolorbox}[breakable,colback=teal!4,colframe=teal!35!black,title={Before you move on}]
What does \bn{খাপ} \bn{খাওয়া} mean? (\textquotedblleft{}To fit; to suit\textquotedblright{}.) Say it once more.
\end{tcolorbox}
